\section{The original realization and closed program-image proof}

\begin{annexitem}{annex:p7.ch88.native-saturation}{thm:p07-ch88-native-saturation}{All closed-program saturation clauses}
\begin{proof}
The set of closed strongly normalizing terms is saturated. The marker
and every marker-headed spine of SN arguments are SN: no head rule
exists, so each step reduces an argument. To verify head expansion,
first note that substitution preserves every designated pure beta
contraction. Hence if \(u[b/x]\) is SN then \(u\) is SN; an
infinite path in the template would survive the substitution. Membership
of its applied spine also makes each spine argument SN by subterm closure.
If the contractum is SN and \(b\) is SN, prove the expanded spine SN
by induction on the sum of reduction heights of the template, argument
and spine terms. A head contraction gives the SN contractum. A step in
the template, substituted argument or spine gives an SN reduct of that
contractum and strictly lowers the relevant height. Substituted copies
of an argument are reduced finitely, possibly zero times. This exhausts
immediate reducts. Finite heights exist by finite branching of finite
syntax. Thus the expanded spine is SN. No arbitrary expansion theorem
or source normalization premise has been invoked.

A nonempty intersection retains every marker spine, is contained in SN
by any one factor, and inherits head expansion in each factor. For the
dependent product, choose \(\Omega\in A\). If \(t\) belongs to
the product then \(t\Omega\) belongs to the saturated \(F(\Omega)\)
and is SN, so its subterm \(t\) is SN. For any \(b\in A\), the
marker and a marker spine with SN arguments followed by \(b\) lie
in \(F(b)\); therefore they lie in the product. For its expansion
clause, append a fixed \(b\in A\) to the given contractum spine.
Membership in the product gives that larger spine in \(F(b)\), and
head expansion in this same \(F(b)\) gives the expanded spine applied
to \(b\). Quantify over \(b\) to recover product membership.
For introduction, \(b\in A\) is SN, so the assumed
\(u[b/x]\in F(b)\) expands to \((\lambda x.u)b\in F(b)\).
This holds for every \(b\), exactly the product definition.
\end{proof}
\end{annexitem}

\begin{annexitem}{annex:p7.ch88.native-values}{thm:p07-ch88-native-values}{Domains, substitution and the two different equality conclusions}
\begin{proof}
The structural source-image certificates used here are the proved
forward validity/typing transport and its structural CC preliminaries,
not CC strong normalization. They classify a type annotation as an
object type or a logical context and classify a value as a proof or
constructor. Constructor/proof argument choices cannot vary between
typing derivations: uniqueness, product injectivity and sort discrimination
fix the domain class. Thus the stripping mask and the structural value
clauses do not depend on an arbitrary choice of derivation. A source
logical context's functionality depends only on its product structure
and annotation classes; object substitution preserves these classes.
It therefore leaves functionality domains unchanged. Typed equality
also preserves them: reflexivity is immediate, symmetry/transitivity
compose; product congruence uses its domain/body inductions; lambda,
application and beta relate objects and both have functionality
\(\mathcal Y\). This supplies the source functionality assertion
without treating object inhabitants as kinds.

Prove total constructor values and their prescribed functionality
membership structurally, for each actual role-correct typing certificate.
A constructor variable is its ambient coordinate. A constructor lambda
curries the body induction over all ambient domain values. If the domain
is an object, the full model requires beta-invariance: the body constructor
interpretation is unchanged when that program coordinate is replaced by
a beta-equal program, as the simultaneous invariance induction below
shows. The restricted model has no such additional requirement. Constructor
application uses its function induction and the actual argument value;
the functionality-substitution identity identifies its result domain.
A small-type product over an object domain has a saturated domain value
and a total saturated family by its proper-subterm inductions. The
dependent-product lemma gives its saturated value. Over a logical domain,
the body gives a saturated value at every semantic input and the domain
is nonempty, so intersection closure gives the value. These are all
constructor cases. A proof value is a finite stripped lambda syntax
with its finitely many external program variables replaced by closed
program coordinates, hence is in \(\mathcal Y\), without an SN claim
for arbitrary ambient valuations.

The needed scoping and substitution assertions are proved simultaneously
on these structural clauses. An indexed lookup selects its older-prefix
coordinate; relocation adds unused younger coordinates, not a new dependency.
A proof-variable substitution is exactly capture-avoiding substitution
of its stripped program in the stripped body. A constructor-variable
substitution changes the semantic constructor coordinate but is absent
from the proof's stripped syntax. The application inductions either
retain two program children or erase the constructor argument, as determined
by the fixed domain class. In a binder, use its old scope on the annotation
and its extended scope on the body. For a constructor lambda this gives
the same function pointwise; for a proof lambda it gives the same raw
program substitution under a fresh pure binder or the erased constructor
binder. For a small-type product it gives the same dependent domain/family
or the same indexed intersection. Applying the identities to each later
environment annotation proves the dependent suffix case as well.

For precision, first prove proof erasure respects full proof equality
as an equality of open pure lambda terms, before closing them with a
valuation. Reflexivity and symmetry/transitivity are direct. Proof lambda
congruence lifts its body equality under a retained pure binder, or
erases its constructor binder. Proof application combines the two pure
child equalities when the argument is a proof and retains only the
function equality when the argument is a constructor. A proof beta
instance becomes a pure beta instance for an object domain and a zero-step
erasure equality for a constructor domain, by the substitution identities.
Products are constructors or logical contexts, not proof-valued endpoints.
These are all interpretable proof cases. Closing the external program
variables preserves this beta equality. No omega rule inferring open
equality from all closed instances is used.

Now prove constructor equality compatibility and, in the full model,
invariance under beta-equal ambient program coordinates, simultaneously
on the constructor syntax and typed equality cases. Constructor variables
keep the same coordinate; proof expressions in constructor arguments
give beta-equal raw values by the preceding erasure lemma. Constructor
lambda gives equal values pointwise, so the full object-domain case is
indeed a beta-invariant function. Constructor application to constructor
arguments uses equal semantic values; to proof arguments the full model
uses its beta-invariant function to identify the beta-equal raw inputs.
Type products have identical saturated domains/families or identical
functionality-indexed families and hence identical sets. Their scopes
use the actual old binder context, transported to the equal annotation;
object-type annotations give equal saturated values, and logical annotations
give equal functionality domains. The corresponding realized extension
families also agree exactly.

In the restricted model, equality cannot relate proof-valued typed endpoints:
its object reflexivity, lambda, application and beta guards require a
logical-context result type. Product equality is not proof-valued either;
symmetry/transitivity preserve this fact. Thus its argument-equality
application case uses equal constructor values, not beta-equal proof inputs
in an arbitrary function. Its function-only case leaves the proof argument
unchanged. Its constructor beta case is the literal semantic substitution
identity. These guards justify the restricted model without imposing the
full model's function restriction. In full equality, constructor beta is
again semantic substitution, while proof beta has only the pure beta
equality just proved. Source conversion relevant to types therefore uses
literal constructor-set equality; it never uses membership invariance
under general pure beta equivalence. The endpoint qualification matters:
an equality of arbitrary ill-typed raw products does not manufacture a
constructor interpretation. The sorted endpoints required by conversion
are supplied by the already proved structural transport. This closes
the value, substitution and invariance inductions without source SN.
\end{proof}
\end{annexitem}

\begin{annexitem}{annex:p7.ch88.native-soundness}{thm:p07-ch88-native-soundness}{All proof-level soundness cases and realized dependencies}
\begin{proof}
Induct on the source proof typing, using the preceding value lemma for
constructor/type subterms. An object-class type has a small-type source-image
sort certificate: type correctness sorts its image, and it cannot have
logical-context sort because the object and logical-context images are
disjoint in the structural kind classification. Thus its native saturated
interpretation is defined even when the source has not separately supplied
the second original premise \(\Gamma\vdash P:\ast\). This uses structural
typing only, not normalization or inverse source adequacy.
A proof variable selects its program coordinate.
The definition of \(D_q\) places that coordinate in the saturated
annotation value at its older prefix; relocation/lookup invariance identifies
this with the type value in the whole environment. Constructor variables
instead range over their functionality domains and are covered by the
value lemma, not by an erroneous proof-variable membership case.

For a proof lambda with object domain \(A\), fix \(e\in D_q(\Gamma)\).
Every \(b\in\rho(A)(e)\) extends this to the realized valuation
\((e,b)\). The body induction gives its stripped instance in
\(\rho(B)(e,b)\). All instances are substitutions in the same pure
one-variable template, by the scoping lemma. Saturated dependent-product
introduction therefore gives the retained pure lambda in
\(\Pi(\rho(A)(e),b\mapsto\rho(B)(e,b))\), exactly the product-type
interpretation. For a lambda over a logical context \(C\), every
\(a\in\varphi_q(C)\) extends the realized valuation. The body program
is independent of that constructor coordinate, since the projection
erases it. The same closed program consequently belongs to all the
body-type values and hence their intersection. No pure lambda is added
in this case, matching source stripping.

For application with object-type argument, the argument induction gives
the actual program \(b\in\rho(A)(e)\). The function induction places
its program in the dependent product, so application belongs to the
family at that very \(b\). Semantic substitution identifies this with
the substituted result type. For application with logical-context argument,
the argument's constructor value belongs to \(\varphi_q(C)\) by the
value lemma. The function program belongs to the intersection over all
those values; choose the actual argument value. Its erased application
is just the unchanged function program, and substitution again identifies
the result type. These are the two application cases, not one ordinary
set-function application at every level.

Product formation and constructor-valued abstraction/application are the
value lemma's cases: they yield saturated types or the appropriate
functionality value, not an additional proof membership premise. In
typing conversion, the proof's stripped program is unchanged and the
old/new type interpretations are literally the same saturated set by
the typed constructor equality lemma. Thus the induction membership
is retained. Context validity and extensions use the stated prefix
dependencies in \(D_q\); logical extensions add any functionality
value, whereas object-type extensions add only a member of its actual
prefix type interpretation. This accounts for every source construction
and conversion case. No reduct-closure of a saturated set, annotation
normalization theorem, or source checker is used.
\end{proof}
\end{annexitem}

\begin{annexitem}{annex:p7.ch88.native-closed-image}{thm:p07-ch88-native-closed-image}{The empty realized environment corollary}
\begin{proof}
The empty realized environment has its unique tuple. Soundness puts
the closed proof image in the interpretation of its small type, which
is saturated by the value lemma. The saturation SN clause gives strong
normalization. In this empty proof-level case the image is precisely
\(\nu_\ast(M)\), not the annotated source term. No reflection from
erased code to source annotation termination is asserted.
\end{proof}
\end{annexitem}

\begin{annexitem}{annex:p7.ch88.native-consistency}{thm:p07-ch88-native-consistency}{The minimal saturated set and the missing marker}
\begin{proof}
Let \(B\) consist of closed SN terms whose beta normal form has
inert head \(\Omega\). It is saturated: the marker spines have
that head, and safe head expansion is SN by the saturation closure
proof and has the same normal form by pure beta confluence. Every
saturated set contains \(B\). Indeed head-reduce a member of \(B\)
until it has a marker head. The process terminates since the term is
SN; a lambda head could never produce a marker-headed normal form.
Its argument subterms are SN, so its final marker spine belongs to
every saturated set. Reverse the finitely many head contractions by
the saturation expansion clause; their arguments are SN subterms.
Hence \(B=\bigcap\mathcal U\).

The type \([A:\ast]A\) is well formed by the source product and
projection rules. Its interpretation is the intersection over all
\(R\in\mathcal U\) of the selected value \(R\), namely \(B\).
A hypothetical closed proof therefore has its stripped image in \(B\)
by soundness. Source syntax and stripping contain no marker constant,
so that pure image is marker-free. Beta substitution cannot introduce
a symbol absent from a term; thus its normal form remains marker-free,
contradicting the required marker head. This proves the original
uninhabited-proposition conclusion using the realization, independently
of the annotation-retaining CC normalization argument.
\end{proof}
\end{annexitem}
